% ============================================================
% Chapter 7: Results
%   Section: RQ2 --- Verifying
% Included from chapters/07-results/results.tex
% Sources: the X-Fact reports in backend/data/evaluation/reports/;
% docs/decisions/evaluation.md (Pilot). The custom-set runs are P0.5 of
% docs/roadmap.md and have not been made.
% ============================================================

\section{RQ2 --- Verifying the Claims}
\label{sec:results-rq2}

RQ2 asks whether the pipeline can verify the claims of the news it
retrieves, and how that can be quantified. The quantification is fixed
in Section~\ref{sec:rq-metrics}: accuracy, macro-$F_1$ and Cohen's
$\kappa$ against the gold labels, natural and synthetic claims apart,
under three evidence conditions (none, retrieved, the annotator's),
with retrieval recall and the stage each error comes from.

\paragraph{What the X-Fact pilot already shows.} The only verification
runs so far are on X-Fact (Section~\ref{sec:results-pilot}), and they
answer RQ2 for political fact-checks rather than for the system's own
domain. Four findings carry over.

\begin{itemize}
    \item \textbf{The zero-shot 3B verifier is weak.} On the 64-claim
    pilot it reached 29.7\% accuracy [18.8, 40.6] and a $\kappa$ of 0.07
    whose interval includes zero; on the 165 Spanish test claims, 19.4\%
    [13.3, 25.5], below always answering the majority class.
    \item \textbf{It leans to \texttt{FALSE}.} It answered
    \texttt{FALSE} to 43 of 64 pilot claims and never
    \texttt{MISLEADING}; without evidence the same model answered
    \texttt{FALSE} to 60 of 64. The lean is the model's prior.
    \item \textbf{Retrieved evidence helps, a little.} Paired over the
    same claims, the pipeline beat the closed-book model by 10.9 accuracy
    points [3.1, 20.3] on the pilot, and by an amount indistinguishable
    from zero on the Spanish test split (3.0 points [$-3.0$, 9.1]).
    \item \textbf{Errors are both retrieval and reasoning.} Under the
    trace rule 44 of 45 wrong verdicts are reasoning errors; under the
    reference rule, 27 of the 29 it can judge are retrieval errors. The
    model misread what it was shown, and what it was shown was not what
    the fact-checker had used --- no annotator link was retrieved at any
    depth.
\end{itemize}

The pilot also showed how much of an open-retrieval score on an old
benchmark is about time rather than verification: 16 of the 20 claims
for which it can be decided cited a source newer than the claim, and
the fact-checker's own verdict page was ranked for 9 of the 24 English
claims (Section~\ref{subsec:pilot-time}). The custom set avoids the
worst of this by construction: its claims are weeks old, not years, and
its annotation rule restricts the gold label to evidence available at
publication time.

\paragraph{On the custom set.} The runs that answer RQ2 for
constructive news --- the test split under the three evidence
conditions, with natural and synthetic claims reported apart, retrieval
recall against the annotator's links, and the stage attribution ---
depend on two P0 items not yet done: the annotator's review of the 108
labels proposed by the AI pass, and the frozen tuning/test split
(Table~\ref{tab:priorities}). A run on unreviewed labels would measure
the pipeline against another model.
\pending{Table: accuracy, macro-$F_1$ and $\kappa$ with intervals on
the custom test split --- closed-book, retrieved and oracle evidence;
natural and synthetic apart; by language and topic group. Retrieval:
link and domain recall at candidates, ranked and cited, and MRR. Stage
attribution under both rules, with the disagreements. Claims set apart
(search or model unreachable) by identifier.}
